\documentclass[10pt]{article}
\usepackage[T1]{fontenc}
\usepackage[utf8]{inputenc}
\usepackage[french]{babel}
\usepackage{amsmath,amssymb,mathrsfs,mathpazo}
\usepackage{geometry}
\geometry{a4paper,margin=2cm}
\usepackage{graphicx,array,tabularx,tabularray,colortbl,multirow,multicol,booktabs,enumitem,pifont,ulem,fancybox,diagbox,cancel}
\usepackage{pstricks,pstricks-add,pst-plot,pst-tree,pst-node,pst-text,pst-func,pst-math,pst-bspline,pst-3dplot,pst-eucl,pst-solides3d}
\let\pscalculate\undefined
\usepackage{pst-bezier}
\usepackage{tikz,pgfplots,tkz-tab}
\usetikzlibrary{arrows,arrows.meta,patterns,calc,positioning,trees,decorations.markings,shapes.geometric,shapes.misc,shapes.arrows,intersections,decorations.pathmorphing,decorations.pathreplacing,backgrounds,shadows,babel}
\pgfplotsset{compat=1.18}
\usepgfplotslibrary{fillbetween}
\usepackage[np]{numprint}
\usepackage{listings,makecell,pas-tableur,twoopt}
\let\barre\undefined
\usepackage[pstricks]{tablvar}
\definecolor{darkgreen}{RGB}{0,100,0}
\providecommand{\red}{\color{red}}
\providecommand{\blue}{\color{blue}}
\providecommand{\green}{\color{green}}
\providecommand{\black}{\color{black}}
\providecommand{\white}{\color{white}}
\providecommand{\gray}{\color{gray}}
\providecommand{\cyan}{\color{cyan}}
\providecommand{\magenta}{\color{magenta}}
\providecommand{\yellow}{\color{yellow}}
\providecommand{\violet}{\color{violet}}
\providecommand{\orange}{\color{orange}}
\providecommand{\euro}{\texteuro}
\providecommand{\EUR}{\texteuro}
\providecommand{\ieme}{\textsuperscript{e}}
\providecommand{\degres}{\ensuremath{^\circ}}
\providecommand{\degree}{\ensuremath{^\circ}}
\providecommand{\vv}[1]{\overrightarrow{#1}}
\providecommand{\Coord}[1]{#1}
\providecommand{\decosix}{}
\providecommand{\starredbullet}{\ensuremath{\bullet}}
\pagestyle{empty}
\setlength{\parindent}{0pt}
\usepackage{adjustbox}\DeclareUnicodeCharacter{2236}{:}\DeclareUnicodeCharacter{2190}{\ensuremath{\leftarrow}}
\providecommand{\R}{\mathbb{R}}
\providecommand{\N}{\mathbb{N}}
\providecommand{\Z}{\mathbb{Z}}
\providecommand{\Q}{\mathbb{Q}}
\providecommand{\D}{\mathbb{D}}
\providecommand{\C}{\mathbb{C}}
\providecommand{\e}{{\mathrm{e}}}
\providecommand{\vect}[1]{\overrightarrow{#1}}
\providecommand{\vectt}[1]{\overrightarrow{\mathrm{#1}}}
\providecommand{\vvt}[1]{\overrightarrow{\mathrm{#1}}}
\providecommand{\np}[2][]{{#2~#1}}
\providecommand{\vv}[1]{{\overrightarrow{#1}}}
\renewcommand{\d}{{\mathrm{d}}}
\providecommand{\degres}{{^\circ}}
\providecommand{\degre}{{^\circ}}
\providecommand{\blue}{\color{blue}}
\providecommand{\red}{\color{red}}
\providecommand{\enskip}{\quad}
\providecommand{\llbracket}{\lbrack\!\lbrack}
\providecommand{\rrbracket}{\rbrack\!\rbrack}
\providecommand{\euro}{{\text{€}}}
\providecommand{\no}{{\text{n°}}}
\providecommand{\interval}[2]{{[#1;#2]}}
\providecommand{\intervFO}[2]{{[#1;#2[}}
\providecommand{\dotfill}{\ldots}

\providecommand{\ds}{\displaystyle}
\providecommand{\textding}[1]{#1}
\providecommand{\Oij}{\(\mathrm{O};\vec{\imath},\vec{\jmath}\)}
\providecommand{\Oijk}{\(\mathrm{O};\vec{\imath},\vec{\jmath},\vec{k}\)}
\providecommand{\Ouv}{\(\mathrm{O};\vec{u},\vec{v}\)}
\providecommand{\intervFF}[2]{\ensuremath{[#1;#2]}}
\renewcommand{\euro}{\ifmmode\text{\texteuro}\else\texteuro\fi}
\DeclareUnicodeCharacter{00B0}{\ifmmode^{\circ}\else\textdegree\fi}
\DeclareUnicodeCharacter{20AC}{\ifmmode\text{\texteuro}\else\texteuro\fi}
% Énoncés : Éducation nationale. Transcription et corrigés : APMEP, Denis Vergès et contributeurs. Découpage et classement : CoopMaths, Sébastien Lozano, Éric Elter et contributeurs. Adaptation de la présentation : Maths Explorer.
% CC BY-SA : https://coopmaths.fr/www/about/
% https://coopmaths.fr/alea/static/e3c/2020/tex/e3c_2020_00_sujet42_2_cor.tex
\begin{document}
\noindent\textit{APMEP / CoopMaths. CC BY-SA. Présentation adaptée par Maths Explorer.}\par\medskip
\subsection*{1.}

Le lait de brebis est représenté par un secteur de \(60^\circ\) : sa proportion est donc de \(\dfrac{60}{360} = \dfrac{1}{6}\) ;

Le lait de chèvre est représenté par un secteur de \(90^\circ\) : sa proportion est donc de \(\dfrac{90}{360} = \dfrac{1}{4}\) ;

Le lait de vache est représenté par un secteur de \(360 - (60 + 90) = 360 - 150 = 210^\circ\) : sa proportion est donc de \(\dfrac{210}{360} = \dfrac{7}{12}\).

On a \( P_C(A) = 0{,}6 \) et \( P(B) = \dfrac{1}{6} \).


\subsection*{2.}

On peut dresser un arbre pondéré de probabilités :

\begin{center}
\begin{tikzpicture}[thick, scale=1.5]
\node (P_-1_0) at (-2,-2.5) {$\phantom{A}$};
\node (P_0_0) at (0,-0.5) {$V$};
\draw (P_-1_0) -- (P_0_0) node[midway, above] {$\frac{7}{12}$};
\node (P_1_0) at (2,-0) {$F$};
\draw (P_0_0) -- (P_1_0) node[midway, above] {$0{,}2$};
\node (P_1_1) at (2,-1) {$A$};
\draw (P_0_0) -- (P_1_1) node[midway, below] {$0{,}8$};
\node (P_0_2) at (0,-2.5) {$C$};
\draw (P_-1_0) -- (P_0_2) node[midway, above] {$\frac14$};
\node (P_1_2) at (2,-2) {$F$};
\draw (P_0_2) -- (P_1_2) node[midway, above] {$0{,}4$};
\node (P_1_3) at (2,-3) {$A$};
\draw (P_0_2) -- (P_1_3) node[midway, below] {$0{,}6$};
\node (P_0_4) at (0,-4.5) {$B$};
\draw (P_-1_0) -- (P_0_4) node[midway, below] {$\frac16$};
\node (P_1_4) at (2,-4) {$F$};
\draw (P_0_4) -- (P_1_4) node[midway, above] {$0{,}7$};
\node (P_1_5) at (2,-5) {$A$};
\draw (P_0_4) -- (P_1_5) node[midway, below] {$0{,}3$};
\end{tikzpicture}
\end{center}

Calculons les probabilités :
\[
P(V \cap A) = P(V) \times P_V(A) = \dfrac{7}{12} \times 0{,}8 = \dfrac{5{,}6}{12} = \dfrac{1{,}4}{3} = \dfrac{7}{15},
\]
\[
P(C \cap A) = P(C) \times P_C(A) = \dfrac{1}{4} \times 0{,}6 = \dfrac{0{,}6}{4} = 0{,}15,
\]
\[
P(B \cap A) = P(B) \times P_B(A) = \dfrac{1}{6} \times 0{,}3 = \dfrac{0{,}3}{6} = 0{,}05.
\]
D'après la loi des probabilités totales :
\begin{align*}
P(A) &= P(V \cap A) + P(C \cap A) + P(B \cap A) \\
&= \dfrac{7}{15} + 0{,}15 + 0{,}05 \\
&= \dfrac{7}{15} + 0{,}20 \\
&= \dfrac{7}{15} + \dfrac{3}{15} \\
&= \dfrac{10}{15} = \dfrac{2}{3}.
\end{align*}

\subsection*{3.}

On a :
\[
P_A(V) = \frac{P(V \cap A)}{P(A)} = \frac{\dfrac{7}{15}}{\dfrac{2}{3}} = \frac{7 \times 3}{15 \times 2} = \frac{21}{30} = \dfrac{7}{10},
\]
soit \(0{,}7\) à \(10^{-3}\) près.

\end{document}
